%%%%

%%%   Самарский государственный технический университет

%%%   Кафедра прикладной математики и информатики

%%%

%%%   Преамбула для статьи в журнал "Вестник Самарского государственного

%%%   технического университета. Серия: «Физико-математические науки»"

%%%   Ориентирована под MikTEx 2.4 for Win32

%%%

%%%   Хотя сам журнал собирается под GNU/Linux на дистрибутиве TeXLive



\documentclass[11pt,a4paper,twoside]{article}



\usepackage[cp1251]{inputenc}

\usepackage[T2A]{fontenc}

\usepackage[english,russian]{babel}

\usepackage{samgtubib}

\usepackage[dvips]{graphicx}

\usepackage[multiple]{footmisc}



\usepackage{samgtu}

\renewcommand{\VestnikYear}{2009} % Год издания Вестника

\renewcommand{\VestnikNumber}{2\,(19)} % Номер Вестника

\renewcommand{\VestnikMonth}{Ноябрь} % Месяц выхода Вестника

\renewcommand{\VestnikData}{01.11.09} % Дата подписания Вестника в печать

\renewcommand{\VestnikUPL}{26,64} % Усл. печ. л.

\renewcommand{\VestnikUIL}{25,73} % Уч.-изд. л.

\renewcommand{\VestnikRegNumber}{479/09} % Регистрационный номер

\renewcommand{\VestnikOrder}{994} % Номер заказа











%\begin{document}

%список ключевых слов

%информационная технология

%поддержка принятия решения

%система управления кадрами

%отбор кадров

%информационная модель

%сетевой график

%метод экспертных оценок

%метод анализа иерархий



\renewcommand{\firstpage}{189} %Начальная страница статьи для выноса в колонтитул

\renewcommand{\lastpage}{191} %Конечная страница статьи для выноса в колонтитул





%\Partition{Краткие сообщения}



%\ShortPartition{Информатика}



\udk{331.108.26:004.9}





\title{Информационные технологии поддержки выработки и принятия решений в системе управления кадрами}

      {Информационные технологии поддержки выработки и принятия решений \dots}

\engtitle{Information technologies for decision support in the staff management system} 

\author{Ю.\,И.~Куц}

       {К\,у\,ц Ю.\,И.}

\engauthor{Yu.\,I.~Kuts}





\date{28.06.2007}





\address{Военно-Морской Институт Радиоэлектроники им. А.\,С.~Попова, \\ г. Санкт-Петербург, Петродворец}

\engaddress{Popov Higher Naval Academy of Radio Electronics, Petrodvorets, Russia}





\email{radio@atnet.ru}





\workabstract{Рассмотрена задача поддержки выработки и принятия решений в системе управления кадрами. Рассмотрены современные информационные технологии поддержки принятия решений. Для оптимизации процесса принятия решений предложена информационная модель на базе сетевого графика. Дано краткое описание метода и~алгоритма обработки информации в~процессе принятия решений.

}



\engabstract{The article investigates the task of decision support in the staff management system. Modern information technologies for decision support are investigated. The information model based on the Net Diagram is suggested purposely optimization of decision support process. Shortly described method and algorithm of data processing in decision support process.}



\renewcommand{\baselinestretch}{.87}



\maketitle



Современные информационные технологии, основанные на использовании средств вычислительной техники и связи внедряются во все виды управленческой деятельности. Процесс создания, внедрения и использования в различных областях управления новых информационных технологий, систем и средств получения, передачи, сбора, обработки, хранения и использования информации называется автоматизацией управления [1].



Целью автоматизации является повышение эффективности функционирования системы управления. При этом основным критерием оценки эффективности управления является время обработки информации в системе [1, 2].



Одним из наиболее сложных и ответственных этапов управления, нуждающихся в информационной поддержке и автоматизации, является этап принятия решений [3, 4]. В качестве вспомогательного средства для обеспечения принятия решений в управлении используются «системы поддержки принятия решений» (Decision Support System "--- DSS) [1].



Автоматизированные системы поддержки принятия решений (АСППР): помогают произвести оценку обстановки (ситуаций), осуществить выбор критериев и оценить их относительную важность; генерируют возможные решения (сценарии действий); осуществляют оценку сценариев (действий, решений) и выбирают лучший; обеспечивают постоянный обмен информацией об обстановке принимаемых решений и помогают согласовать групповые решения; моделируют принимаемые решения (в~тех случаях, когда это возможно); осуществляют динамический компьютерный анализ возможных последствий принимаемых решений; производят сбор данных о результатах реализации принятых решений и осуществляют оценку результатов [1].



В системе управления кадрами решения о назначении на должности принимаются должностными лицами (лицами, принимающими решения "--- ЛПР) в соответствии с предоставленными им полномочиями. Отбор кандидатов для назначения на должности и подготовку решений должностных лиц осуществляют кадровые органы.



Процесс отбора кандидатов для назначения на должности включает, как правило, два основных этапа [4]: формирование резерва кандидатов для назначения на должности и конкурсный отбор кандидатов, ранее включённых в резерв.



На этапе формирования резерва определяются требования, предъявляемые к кандидатам для назначения на данные должности, и составляется список кандидатов.



На этапе конкурсного отбора производится сравнительная оценка кандидатов и определяется кандидат, в наибольшей степени соответствующий требованиям к данной должности.



На обоих этапах в системе управления кадрами циркулирует большой объем информации по кадрам, осуществляется ее сбор, регистрация, накопление, хранение, анализ и переработка.



Специфика выработки и принятия решения по кадрам заключается в том, что факторы (критерии), определяющие конечные результаты, в основном не поддаются непосредственному измерению. Таким образом, принятие решений по отбору кандидатов для назначения на должности относится к классу сложных слабо структурированных задач. 



На практике, выработка и принятие решения по кадрам в значительной степени зависят от профессиональной подготовленности, личного опыта и интуиции ЛПР и работников кадрового органа. При этом возникают противоречия:



-- между потребностью увеличения времени для выработки обоснованного решения и необходимостью повышения оперативности управления;



-- между постоянно увеличивающимся потоком информации и сокращением времени на её обработку;



-- между сокращением времени на подготовку решения и возрастающей ответственностью за его принятие.



Для разрешения указанных противоречий необходимо, с одной стороны, увеличить скорость обработки информации в~системе, с~другой стороны, "--- обеспечить высокое качество принимаемого решения.



Существенное повышение оперативности и~качества управления достигается применением в~системе управления кадрами специализированной АСППР.



Создание АСППР для системы управления кадрами является сложной научно"=технической задачей, решение которой предполагает: анализ потоков информации в~процессе принятия решений и~разработку предложений по их совершенствованию; разработку (выбор) метода обработки информации; алгоритмизацию процесса обработки информации; разработку специального математического и программного обеспечения и др.



В результате проведённого автором исследования процесса выработки и~принятия решения в~системе управления кадрами были получены следующие результаты.



1. Разработана и проанализирована информационная модель процесса выработки и принятия решения в существующей системе управления кадрами (модель «как есть»).



2. Синтезирована информационная модель процесса выработки и принятия решения в перспективной системе управления кадрами (модель «как должно быть» на базе новых рациональных технологий обработки информации).



3. Предложен метод анализа иерархий для обработки информации в процессе принятия решения по отбору кадров.



4. Разработан алгоритм обработки информации, позволяющий реализовать указанный метод для решения задачи по отбору кадров с использованием ЭВМ.



Информационные модели, отражающие логическую схему получения, обработки, накопления и~передачи данных в~процессе принятия решений и~обработки информации, разработаны с~целью получения материала, необходимого для совершенствования процесса принятия решений в системе управления военными кадрами.



Модель «как есть» "--- это «снимок» фактического состояния информационной системы на момент обследования, позволяющий дать формальное описание потоков информации (главным образом документопотоков) с позиции системного анализа, а также выявить ошибки и «узкие места», сформулировать предложения по улучшению ситуации.



Модель «как должно быть» "--- интегрирует перспективные предложения по оптимизации информационных потоков, позволяет сформулировать видение новых рациональных технологий обработки информации.



Информационные модели построены на основе сетевого графика [2, 5]. Исследование информационного потока на основе его сетевой модели позволило достичь наглядности функционирования системы управления и~движения потоков информации; применение математического аппарата теории графов позволило оптимизировать процесс управления; получена возможность также представить динамику управления и~движения информации, которая ускользает при пользовании другими методами.



Оптимизированная информационная модель процесса принятия решений позволила осуществить выбор соответствующих методов получения и обработки информации.



Для получения исходной информации предложено использование метода экспертных оценок, сущность которого заключается в~проведении экспертами интуитивно-логического анализа с количественной оценкой суждений~[6].



Информация, полученная от экспертов, поддаётся формальной обработке средствами вычислительной техники. При этом предложено использование метода анализа иерархий, состоящего в декомпозиции проблемы на все более простые составляющие части и~дальнейшей обработке последовательности суждений лица, принимающего решение, по парным сравнениям [6--8].



Решение задачи отбора кандидатов для назначения на должность методом анализа иерархий представляет по существу получение численной оценки предпочтения альтернативных вариантов (кандидатов) путём обработки информации, полученной от экспертов и~выбор объекта (кандидата) с~наибольшей оценкой.



Данная задача представляется в иерархической форме тремя уровнями.



На первом (высшем) уровне находится общая цель "--- «кандидат».



На втором уровне находятся факторы (критерии), уточняющие цель.



На третьем (нижнем) уровне находятся сравниваемые кандидаты (альтернативы), которые должны быть оценены по отношению к факторам (критериям) второго уровня.



Исходной информацией для решения задачи является:



 -- список сравниваемых кандидатов;



-- перечень критериев, с учётом которых проводится оценивание кандидатов;



-- шкала интенсивностей относительной важности критериев;



-- шкала значений оценок;



-- численные значения оценок объектов по учитываемым критериям.



Исходные данные обрабатываются в соответствии с разработанным алгоритмом.



В соответствии с методом анализа иерархий алгоритм решения задачи включает: ввод исходных данных; формирование и заполнение матриц попарных сравнений для уровня~2; формирование и заполнение матриц попарных сравнений для уровня~3; синтез приоритетов, наибольшего собственного значения матрицы суждений, индекса согласованности и отношения согласованности для уровня 2; синтез приоритетов, наибольшего собственного значения матрицы суждений, индекса согласованности и отношения согласованности для уровня 3; синтез глобальных приоритетов; выбор объекта, имеющего наибольший приоритет.



Разработанный на языке математических описаний алгоритм может быть применён для решения задачи с использованием ЭВМ. Для функционирования алгоритма на ЭВМ необходимо записать его на алгоритмическом языке из числа так называемых языков программирования. Выбор языка программирования осуществляется с учётом особенностей ЭВМ и~функционирующей на ней операционной системы~(ОС).



В заключение следует отметить, что полученные в настоящей работе результаты являются новыми. Впервые в~данной предметной области использованы: метод сетевого планирования для моделирования информационных потоков в~процессе выработки и~принятия решения по кадрам; метод анализа иерархий для решения задачи по отбору кадров.



Указанные методы апробированы и~успешно применяются в~других предметных областях, что подтверждает достоверность полученных результатов.



Разработанные для процесса принятия решений в~системе управления кадрами информационная модель, метод и~алгоритм обработки информации в~комплексе с~необходимыми техническими средствами являются основой для создания АСППР управления кадрами, внедрение которой позволит повысить эффективность функционирования системы управления.



\vspace{-3mm}



\begin{thebibliography}{}

\item{}Автоматизация управления и~связь в~ВМФ [Teкст]~/ Под общ. ред. Ю.\,М.~Кононова; изд. 2-е. "--- СПб.: Элмор, 2001.



\item\emph{Модин, А.\,А.} Справочник разработчика АСУ [Teкст]~/ А.\,А.~Модин, Е.\,Г.~Яковенков, Е.\,П.~Погребной. "--- М.: Экономика, 1978. "--- 583~c.



\item\emph{Волынец, Ю.\,Ф.} Системный анализ процессов кадрового обеспечения в едином информационно"=функциональном пространстве ВМФ [Текст]~/ Ю.\,Ф.~Волынец,  А.\,М.~Михальчук~ /\!/ Военная радиоэлектроника: опыт использования и~проблемы, подготовка специалистов. Материалы двенадцатой научно"=технической конференции (межвузовской). "--- Петродворец: ВМИРЭ, 2001. "---  Ч.~1. "--- С.~137.



\item\emph{Днов, В.\,Н.} Представления знаний в~системе интеллектуальной поддержки кадрового органа [Текст]~/ 

В.\,Н.~Днов~/\!/ Военная радиоэлектроника: опыт использования и проблемы, подготовка специалистов. Материалы шестнадцатой научно"=технической конференции (межвузовской). "--- Петродворец: ВМИРЭ, 2005. "--- С.~231--232.



\item\emph{Абчук, В.\,А.} Справочник по исследованию операций [Текст]~/ 

В.\,А.~Абчук, Ф.\,А.~Матвейчук, Л.\,П.~Томашевский; Под общ. ред. Ф.\,А. Матвейчука. "--- М.: Воен. издат., 1979. "--- 368~c.



\item\emph{Мухин, В.\,И.} Исследование систем управления. Анализ и синтез систем управления: учебник для вузов [Текст]~/ В.\,И. Мухин; Изд. 2-е, доп. и перераб. "--- М.: Экзамен, 2006. "--- 479~с. "--- ISBN 5--472--01927--3.



\item\emph{Ларичев, О.\,И.} Качественные методы принятия решений [Текст]~/ О.\,И.~Ларичев, Е.\,М.~Мошкович.

"--- М.: Физматлит, 1996. "--- 208~с.



\item\emph{Саати, Т.\,Г.} Принятие решений. Метод анализа иерархий [Текст]~/  Т.\,Г.~Саати; пер. с англ. Р.\,Г.~Вачнадзе. "--- М.: Радио и связь, 1993. "--- 320 с.

\end{thebibliography}



\vspace{-4mm}

\lastdate{08.01.2008}



\vspace{1mm}



\makeengtitle



\renewcommand{\baselinestretch}{.9}



%\end{document}







